% CVPR 2027 submission — SUPPLEMENTARY MATERIAL (separate PDF); template: https://github.com/cvpr-org/author-kit (CVPR2026-v1 kit, year bumped)
\documentclass[10pt,twocolumn,letterpaper]{article}

%%%%%%%%% PAPER TYPE
% \usepackage{cvpr}              % camera-ready
\usepackage[review]{cvpr}      % review version
% \usepackage[pagenumbers]{cvpr} % arXiv version

\usepackage{booktabs}
\usepackage{graphicx}
\usepackage{multirow}
\usepackage[table]{xcolor}
\usepackage{amsmath}
\usepackage{url}
\usepackage{tikz}
\usetikzlibrary{positioning,arrows.meta,patterns}
\usepackage{cuted}   % two-column teaser (kit preamble.tex)
\usepackage{caption}
\usepackage{microtype}  % kit preamble.tex: tighter layout
\definecolor{ourteal}{HTML}{0D5C63}
\definecolor{memgreen}{HTML}{2F8F83}
\definecolor{lexorange}{HTML}{C4841D}

% Numbers still to be re-read from regenerated tables (analyze_results.py). Remove before submission.
\newcommand{\vnum}[1]{\textcolor{red}{#1}}
\newcommand{\todo}[1]{\textcolor{blue}{[TODO: #1]}}

\definecolor{cvprblue}{rgb}{0.21,0.49,0.74}
\usepackage[pagebackref,breaklinks,colorlinks,allcolors=cvprblue]{hyperref}

%%%%%%%%% PAPER ID
\def\paperID{*****}
\def\confName{CVPR}
\def\confYear{2027}

\title{Diagnosis-Guided Compressed Latent Replay for Continual Document Understanding}

\author{Anonymous CVPR submission}

\begin{document}
\maketitlesupplementary
\appendix

\section{Setup details}
\label{app:setup}

\paragraph{DIL unified schema.} A single four-entity label space \{HEADER, KEY, VALUE, OTHER\} (9 BIO tags) is held fixed over FUNSD$\rightarrow$SROIE$\rightarrow$CORD. FUNSD: HEADER$\rightarrow$HEADER, QUESTION$\rightarrow$KEY, ANSWER$\rightarrow$VALUE. SROIE: COMPANY$\rightarrow$KEY; ADDRESS/DATE/TOTAL$\rightarrow$VALUE. CORD: menu/sub\_total/total$\rightarrow$VALUE; void\_menu/sub$\rightarrow$OTHER. KEY originates only from FUNSD and SROIE, so the schema is VALUE-dominant; we report per-class F1 for every DIL result and treat a high aggregate as retention only if the sparse classes are retained. The reverse order is used as a robustness check.

\paragraph{CIL-CORD sessions.} CORD's 30 fine-grained classes in five sessions of six, ordered easy-to-hard by super-class: (0) core menu fields; (1) extended menu + sub; (2) sub\_total fields; (3) total / payment methods; (4) residual total + void/auxiliary. Every receipt containing at least one in-session entity is kept, with out-of-session tokens relabelled to ``O'', giving per-session training pools of 800/382/549/723/780 receipts. Class imbalance across sessions is deliberately left unmodified. Joint uses the full-label pool (each receipt once with all its labels) rather than the concatenation of the masked per-session copies, which would feed conflicting labels for the same document.

\paragraph{Mixed.} Six sessions: CORD-menu $\rightarrow$ FUNSD $\rightarrow$ CORD-menu-extended $\rightarrow$ SROIE $\rightarrow$ CORD-sub\_total $\rightarrow$ FUNSD-revisit; boundaries alternate new-class and new-domain transitions and the final session is an explicit retention check.

\paragraph{Datasets.} FUNSD (\texttt{nielsr/funsd-layoutlmv3}; 149/50), CORD (\texttt{naver-clova-ix/cord-v2}; 800/100/100), SROIE (\texttt{mp-02/sroie}; 626/347). Dataset-provided OCR and boxes throughout; images 224$\times$224; sequence length 512.

\paragraph{Hyperparameters.} AdamW, lr $5\times10^{-5}$ constant, weight decay $0.01$, clipping $1.0$, batch size 8 on the primary grid (1--2 with gradient checkpointing on the 6\,GB secondary-backbone machine; replay batch 2, or 1 for DER++ on BROS, whose HF implementation lacks checkpointing). EWC $\lambda{=}10^3$, Fisher over 200 documents; LwF $T{=}2$, $\alpha{=}1$; ER/DER++ buffer 200 documents, DER++ $\alpha{=}\beta{=}0.5$. Early stopping on task val F1, patience 2, $\delta{=}0.1$, cap 100 epochs (never binding; tasks converge in 4--8). Seeds 42/7/123. Primary-grid document runs used bf16 autocast on a 24\,GB GPU; the secondary-backbone cells ran in fp32 or fp16 on a 6\,GB GPU. Image CIL (CIFAR-100, ImageNet-R; ViT-B/16, 10 sessions): every reported cell, including the frozen-backbone and replay baselines, was trained on a 6\,GB GPU at batch size 16 in fp32 with gradient checkpointing, 100-epoch cap with the same early stopping, seeds 42/7/123; the new CODA-Prompt, CL-LoRA and LCA image cells use this recipe unchanged.

\paragraph{Compute.} Every run fits a single GPU; the full generality grid is roughly 200 GPU-hours on a 24\,GB card.

\section{Additional results}
\label{app:extra}
\paragraph{Prompt-, LoRA- and flatness-based baselines.} \cref{tab:prompt_lora} reports the frozen-backbone prompt family (L2P, DualPrompt, CODA-Prompt), the orthogonal/continual LoRA family and ER with C-Flat++ sharpness-aware optimisation on DIL, three seeds, LayoutLMv3 and BERT. The prompt family forgets little but never acquires the task (AA $\approx30$--$35$, below the naive floor of $41.3$): a frozen document encoder with a drifting shared head cannot be routed back into place. LoRA adapters acquire but forget as much as naive training does, because the adapter does not protect the head. Only the replay-based variant reaches the replay ceiling, and its sharpness term adds nothing over plain ER ($88.4$ vs.\ $87.9$ on LayoutLMv3), consistent with Sec.~5 of the main paper: the remedy ordering is set by whether the head is re-grounded, not by the optimiser.

\begin{table}[t]
\centering\small
\caption{Prompt/LoRA/C-Flat++ baselines, DIL, AA and BWT, mean$\pm$std over seeds 42/7/123. Naive floor $41.3$ (LayoutLMv3), ER $87.9$.}
\label{tab:prompt_lora}
\setlength{\tabcolsep}{2pt}\scriptsize
\begin{tabular}{lcccc}
\toprule
 & \multicolumn{2}{c}{LayoutLMv3} & \multicolumn{2}{c}{BERT} \\
Method & AA & BWT & AA & BWT \\
\midrule
L2P & $30.4\pm0.4$ & $-11.5\pm0.1$ & $34.4\pm0.7$ & $-38.5\pm1.0$ \\
DualPrompt & $29.9\pm0.6$ & $-11.1\pm0.2$ & $34.0\pm0.2$ & $-40.2\pm1.1$ \\
CODA-Prompt & $32.5\pm1.2$ & $-10.5\pm0.2$ & $35.2\pm0.7$ & $-40.7\pm1.3$ \\
O-LoRA & $40.8\pm0.3$ & $-56.8\pm1.0$ & $37.4\pm0.3$ & $-53.4\pm0.9$ \\
CL-LoRA & $41.1\pm0.3$ & $-58.8\pm2.1$ & $38.6\pm0.4$ & $-31.6\pm1.9$ \\
ER + C-Flat++ & $88.4\pm0.1$ & $-1.8\pm0.6$ & $78.3\pm0.4$ & $-2.2\pm0.1$ \\
\bottomrule
\end{tabular}
\end{table}

\paragraph{Image CIL: forgetting and the dominated rows.} \cref{tab:image_supp} extends Table~4 of the main paper with BWT and with every image run outside the frozen reference set that we keep as evidence, all under the same recipe (ViT-B/16, ten sessions, batch 16, early stopping, seeds 42/7/123; cells not yet on disk are marked pending, never imputed). Two readings matter. Replay at $200$ exemplars is far below the $2{,}000$-exemplar rows ($42.3/21.6$ vs.\ $76.2/54.8$ for ER), so the buffer size in Table~4 is the one at which replay is competitive, not a favourable choice for our rows. SLCA's whole margin comes from its classifier alignment: without it the same recipe falls to $38.0/25.8$, below $200$-exemplar replay on CIFAR-100, which is the image-side instance of the head-locus finding. Per-class Gaussian latent replay, the image port of the document-side Gaussian store, sits at the naive floor ($10.7$).

\begin{table}[t]
\centering\small
\caption{Image class-incremental learning, AA and BWT, mean$\pm$std over seeds 42/7/123 (ViT-B/16, ten sessions). Frozen reference set first, our rows, the joint bound, then the dominated rows kept as evidence. SimpleCIL is a single seed.}
\label{tab:image_supp}
\setlength{\tabcolsep}{2pt}\scriptsize
\resizebox{\linewidth}{!}{% Auto-generated by scripts/rca/make_image_table.py (do not edit by hand).
% AA and BWT, mean $\pm$ std over seeds 42/7/123; frozen reference set first, then the extra rows kept as evidence.
\begin{tabular}{lcccc}
\toprule
 & \multicolumn{2}{c}{\textbf{CIFAR-100}} & \multicolumn{2}{c}{\textbf{ImageNet-R}} \\
\textbf{Method} & AA & BWT & AA & BWT \\
\midrule
Naive (lower bound) & 11.7\;{\scriptsize $\pm$ 1.1} & -95.6\;{\scriptsize $\pm$ 1.4} & 9.3\;{\scriptsize $\pm$ 0.3} & -87.1\;{\scriptsize $\pm$ 0.8} \\
LwF & 15.7\;{\scriptsize $\pm$ 0.6} & -91.0\;{\scriptsize $\pm$ 0.7} & 12.7\;{\scriptsize $\pm$ 1.4} & -84.4\;{\scriptsize $\pm$ 1.4} \\
ER (2{,}000 exemplars) & 76.2\;{\scriptsize $\pm$ 1.2} & -23.8\;{\scriptsize $\pm$ 1.3} & 54.8\;{\scriptsize $\pm$ 0.7} & -37.9\;{\scriptsize $\pm$ 0.9} \\
DER++ (2{,}000 exemplars) & 76.5\;{\scriptsize $\pm$ 1.5} & -23.6\;{\scriptsize $\pm$ 1.7} & 52.8\;{\scriptsize $\pm$ 1.1} & -40.5\;{\scriptsize $\pm$ 1.5} \\
CODA-Prompt & \emph{pending} & -- & 8.4\;{\scriptsize $\pm$ 0.7} & -67.4\;{\scriptsize $\pm$ 1.7} \\
CL-LoRA (2025) & \emph{pending} & -- & \emph{pending} & -- \\
SLCA (2023) & 87.5\;{\scriptsize $\pm$ 0.2} & -7.8\;{\scriptsize $\pm$ 0.2} & 65.0\;{\scriptsize $\pm$ 0.3} & -13.5\;{\scriptsize $\pm$ 0.2} \\
RanPAC (2023) & 85.3\;{\scriptsize $\pm$ 0.1} & -5.0\;{\scriptsize $\pm$ 0.1} & 58.5\;{\scriptsize $\pm$ 0.3} & -7.0\;{\scriptsize $\pm$ 0.3} \\
LCA (2026) & 77.2\;{\scriptsize $\pm$ 3.2} & -15.8\;{\scriptsize $\pm$ 3.8} & 50.6\;{\scriptsize (1 seed)} & -27.0\;{\scriptsize (1 seed)} \\
\midrule
\textbf{Latent bank, 500/class (ours, uncompressed)} & 80.0\;{\scriptsize $\pm$ 0.7} & -19.8\;{\scriptsize $\pm$ 0.7} & 55.1\;{\scriptsize $\pm$ 1.3} & -37.4\;{\scriptsize $\pm$ 1.3} \\
\textbf{CoLaR $r{=}128$, 500/class (ours)} & -- & -- & 55.3\;{\scriptsize $\pm$ 1.4} & -36.8\;{\scriptsize $\pm$ 1.5} \\
\midrule
Joint (upper bound) & 89.5\;{\scriptsize $\pm$ 0.4} & -- & 77.6\;{\scriptsize $\pm$ 0.5} & -- \\
\midrule
EWC & 13.9\;{\scriptsize $\pm$ 1.5} & -93.1\;{\scriptsize $\pm$ 1.6} & 9.3\;{\scriptsize $\pm$ 0.2} & -87.2\;{\scriptsize $\pm$ 0.1} \\
ER (200 exemplars) & 42.3\;{\scriptsize $\pm$ 2.7} & -61.6\;{\scriptsize $\pm$ 3.1} & 21.6\;{\scriptsize $\pm$ 1.3} & -75.2\;{\scriptsize $\pm$ 1.4} \\
DER++ (200 exemplars) & 43.6\;{\scriptsize $\pm$ 4.0} & -60.6\;{\scriptsize $\pm$ 4.2} & 20.2\;{\scriptsize $\pm$ 2.5} & -77.2\;{\scriptsize $\pm$ 2.5} \\
SLCA without classifier alignment & 38.0\;{\scriptsize $\pm$ 3.9} & -66.7\;{\scriptsize $\pm$ 4.3} & 25.8\;{\scriptsize $\pm$ 0.7} & -67.4\;{\scriptsize $\pm$ 0.9} \\
SimpleCIL (frozen trunk, NCM) & 81.7\;{\scriptsize (1 seed)} & -5.5\;{\scriptsize (1 seed)} & 48.2\;{\scriptsize (1 seed)} & -6.2\;{\scriptsize (1 seed)} \\
Per-class Gaussian latent replay (falsified) & 10.7\;{\scriptsize $\pm$ 0.8} & -96.5\;{\scriptsize $\pm$ 1.0} & -- & -- \\
\bottomrule
\end{tabular}
}
\end{table}

\paragraph{Ridge decomposition of the plastic-head gap.} \cref{tab:probe} reports the Pass~2b probe on the trained trunk of our replay rows (seed 42). After the last session the trunk is frozen and a closed-form ridge readout (random projection to $10{,}000$ ReLU features, $\lambda{=}10$, the RanPAC recipe) is fitted under three data conditions: on the stored bank only (deployable), on all training data of all sessions (an oracle ceiling on what the features support), and on all data through the frozen pre-trained trunk (image only). Three differences follow: the \emph{head gap} (full ridge minus the trained head) is the total headroom of the readout; the \emph{coverage gap} (full minus stored) is the part that needs features of inputs the bank does not hold; the \emph{trunk gap} (adapted minus frozen) is what replay-driven adaptation bought. On images the head gap is $9.5/16.1$ points and is mostly coverage ($7.5/14.3$), so a better analytic head on the stored latents does not close the gap to SLCA; the adapted trunk is nevertheless $2.9/15.0$ points better than the frozen one under the same readout, so adaptation is not the problem. On documents the head gap of CoLaR is within one point on DIL and Mixed and $7.7$ points on CIL-CORD with zero coverage gap, i.e.\ the compressed store already covers its documents and the residual is the plastic head itself. The probe reads every past training image once at the end and is therefore a diagnostic, not a continual method.

\begin{table}[t]
\centering\small
\caption{Pass~2b ridge decomposition, seed 42, AA. \emph{Head}: the trained plastic head. \emph{Ridge, stored/full/frozen}: closed-form readout on the stored latents, on all seen training data, and on all data through the frozen pre-trained trunk. Gaps as defined in the text; a dash means the condition does not apply (documents have no frozen-trunk reference).}
\label{tab:probe}
\setlength{\tabcolsep}{2pt}\scriptsize
\resizebox{\linewidth}{!}{% Auto-generated by scripts/rca/make_image_table.py from results/*/probes.json (Pass 2b, seed 42; do not edit by hand).
\begin{tabular}{lccccccc}
\toprule
\textbf{Cell} & Head & Ridge, stored & Ridge, full & Ridge, frozen & Head gap & Trunk gap & Coverage gap \\
\midrule
CIFAR-100, latent bank & 78.8 & 80.8 & 88.3 & 85.4 & 9.5 & 2.9 & 7.5 \\
ImageNet-R, latent bank & 56.2 & 58.0 & 72.3 & 57.3 & 16.1 & 15.0 & 14.3 \\
DIL, CoLaR & 90.9 & 91.1 & 91.3 & -- & 0.3 & -- & 0.2 \\
CIL-CORD, CoLaR & 67.5 & 75.2 & 75.2 & -- & 7.7 & -- & 0.0 \\
Mixed, CoLaR & 70.9 & 71.0 & 71.6 & -- & 0.7 & -- & 0.6 \\
\bottomrule
\end{tabular}
}
\end{table}

\paragraph{Per-backbone forgetting.} The full BWT and AF tables for every backbone$\times$scenario$\times$method cell of the generality grid are generated from the released run directories (\texttt{results/table\_backbone\_\{BWT,AF\}.tex}) and will ship with the code; Table~1 of the main paper already carries the AA view.

\paragraph{Budget-adaptive allocation is a frontier, not a gain.} Allocating a fixed compressed-replay budget unevenly across tasks (documents $[18,5,2]$, ranks $[64,128,64]$ vs.\ a uniform $10$ documents at rank $64$; 5 epochs, $k{=}4$) gave $+4.0$ AA on seed 42 while passing a pre-registered per-domain guard. On seeds 7 and 123 the same comparison gave $-8.0$ and $+5.7$: $+0.6\pm7.5$ over three seeds, a null. The seed-42 result was a point on a Pareto frontier (old-domain retention traded against the final domain at $\approx9{:}1$), not a shift of the frontier, and we report it as such.

%% FROZEN BASELINE SET (2026-10-21): CPFD and IS3 are the document-side paired-slot rows; CODA-Prompt is row 6.
%% Their CIL-CORD (and CPFD's Mixed) cells are reported as nulls with this disclosure. Evidence:
%% results/phase1/cil_cord_cpfd_seed42 (+adaptive_ce_min=0.5), results/phase1/cil_cord_cpfd_dw0_seed42
%% (distill_weight=0), results/phase1/cil_cord_coda_prompt_seed42 (n_components=50); ledger cells
%% results/{cil_cord,mixed}_{cpfd,is3,coda_prompt}_seed*/metrics.json. DIL cells for CPFD/IS3 are added in Phase 3.
\paragraph{Null cells among the fixed baselines (CPFD, IS3, CODA-Prompt).} Three of the twelve reference rows do not converge on the document scenarios, and we report those cells as nulls rather than as scores. CPFD and IS3 were ported onto the single growing token-classification head used by every other method here (the reference implementations keep per-task classifiers). On CIL-CORD, CPFD learns the first session ($90$--$95$ F1) and then reads $0.00$ on every session from the first class-adding boundary (AA $2.7$ and $0.0$ on seeds 42 and 123); on Mixed it reads $6.0$, $0.0$ and $0.1$; and on DIL, where no class is ever added, the first domain falls to $8.0/4.0/0.9$ F1 at the first boundary and the final row is zero on every domain for seeds 42 and 7 (AA $0.0/4.2/14.2$, reported as $6.2\pm7.2$ since one seed does learn the last domain), so the collapse is not a property of head growth. Two single-knob ablations on seed 42 reproduce the collapse exactly: flooring the per-sample CE weight at $0.5$ (session-1 F1 $0.00$) and switching the distillation term off entirely (session-1 F1 $2.1$), so neither the adaptive CE factor nor the attention/logit distillation is the cause; what remains is the pseudo-labelling and O-token masking path, which removes the training signal on the new session's entities. IS3 learns each new session (diagonal $93/29/32/66/10$ on seed 123) but never predicts a previously learned class again---every off-diagonal entry is exactly $0.00$---giving AA $4.4\pm1.7$ on CIL-CORD; on Mixed, whose last session adds no classes, it recovers to $30.1\pm1.0$ and that cell is reported. On DIL the nine-tag head never grows, so IS3's old/new split is degenerate (every entity row is ``old'' at every boundary and the O-token distillation spans the whole output) and the method reduces to sequential fine-tuning under a full-output distillation term; seed 42 learns the first domain ($88.2$) and drops both earlier domains below $8$ F1 at the next boundary (AA $38.3$ vs.\ naive $41.3$). We report that single seed rather than spend further compute on a structural limitation of the port. CODA-Prompt reads AA $0.0$ on all three CIL-CORD seeds with its default pool of $10$ components; enlarging the pool to $50$ ($10$ per session) lets some sessions be learned (diagonal $79.5/40.6/0.0/18.9/36.8$) but the final row is $[0,0,0,0,36.8]$, AA $7.4$, with one session never learned, so the cell stays a null. These are properties of the methods' O-tag and static-pool assumptions under five class-adding sessions, not of the sessions themselves: the same sessions are learned by naive training (AA $19.0$) and by the joint upper bound ($33.3$).

\end{document}
